\section{The Dragon Path}

In \emph{The Symbolism of the Cross}, \textbf{Rene Guenon} shows the influence of \textbf{Matgioi} in his expression of Tradition. In particular, the teachings of the Will of Heaven, the degrees of Existence, and the spiral path upward through the degrees, are taken from him.

Matgioi claimed to have been one of just five Europeans initiated into Chinese esoterism (`Society of Heaven and Earth'). He was motivated to describe those teachings according to the following plan:

\begin{quotex}
I was led to divide this work into three parts: one relates, under the title of ``\emph{Metaphysical Path}``, the principles of the Tradition and its philosophical and cosmogonic movement: the second, under the title of ``\emph{Rational Path}'' will relate the systematization of Tradition to Taoism, or Lao Tzu's ``\emph{Tao Te Ching}''; the third, under the title of ``\emph{Social Path}``, will relate the adaptation of the Tradition, with the political philosophy of Confucius.
\end{quotex}


He never published the third volume. \emph{The Metaphysical Path} is based on the \emph{I Ching}, or \emph{Book of Changes}. It is usually described today as a guide to divination. However, that indicates the degradation of a teaching. Just as the Tarot in the West concealed a complete metaphysical system through its symbols, in our time it has also become a tool for fortune telling.

Likewise, the \emph{Emperor Fohi} collected the metaphysical and cosmogonic mysteries of the Primordial Tradition under the title I Ching. The set of 64 symbols contains the mysteries and describes how to reach higher stages. Symbols force one to surpass the intellectual function with its talk about God, in order to ``see'' God. In this selection from \emph{The Metaphysical Path}, we see that path that the ``gifted man'' must follow.

Let us see, for example, in a few lines, how the initiate finds here rules for his conduct as a magician, for his special asceticism.

\emph{Hidden Dragon}. --- The gifted man must regulate his conduct according to the activity of heaven; the gifted man not having been instructed yet, the will of heaven is hidden from his insufficient eye: he therefore remains enveloped in his gangue of the imperfect mortal. The gifted man must therefore meditate, be silent, and try to develop himself in study and contemplation. If he acted while the dragon is hidden, he would not give his measure, and would fall into an error which would be prejudicial to his future.

\emph{Dragon in the Rice Field}. --- The gifted man is aware of his virtue, but cannot yet leave the earth. Little by little he improves beings through his teaching; but he is not yet permitted either to command or to manifest himself. He must only endeavor to follow the fortune and the example of the Mages who preceded him.

\emph{Visible Dragon}. --- The gifted man, placed in a situation inferior to his merits, is in danger; he must act with circumspection; for he attracts by his virtue the sympathy of the universe, and, by this sympathy, the hatred of his superiors. But whether he withdraws or stays, let him always take care to follow the normal path (Tao).

\emph{Jumping Dragon}. --- When the gifted man acts, it is never unrelated to the moment when he acts. He therefore increased his merits and his virtue to be distinguished at a precise and determined moment; he is free to move forward or backward; he retained all his freedom; he can edify by dazzling virtue, just as he can descend again in meritorious humility; in this situation, he must take inspiration from the circumstances.

\emph{Flying Dragon}. --- The gifted man occupies the superior position which suits him. Having arrived at the heights of intelligence, it is sweet to look down on a man equally endowed with virtue, to help him with his examples and to associate him with his power. When one is in the fullness of one's means, one must act.

\emph{Hovering Dragon}. --- Infinite beauty is difficult to maintain. Also the gifted man must know how to advance and retreat in time so as never to expose himself to losing it. One should never be excessive in one's actions, even good ones.

To end a presentation which could go on indefinitely, let us give, without commentary, the six short, simple and rich apothegms by which Confucius, with his usual clearness and conciseness, determines, on the march of the Dragons, the normal conduct of the ordinary citizen. This quote will give a perfect idea of how the Chinese sages understand the moral law.

\begin{enumerate}


\item Do not change according to the age; do not cling to fame; flee the world; have no sorrow at not being appreciated or known by men.

\item Good faith in the smallest words; circumspection in actions; be on guard against lying; improve your society by transforming virtue without boasting about it.

\item Occupy a high position without being proud of it; occupy an inferior position without complaining about it.

\item Perfect your skills; take advantage of the opportune moment.

\item Act, and, by one's action, save the universe.

\item Be careful not to be too noble to have an occupation, and to be too elevated to have friends.
\end{enumerate}



\flrightit{Posted on 2022-08-03 by Cologero}

\begin{center}* * *\end{center}

\begin{footnotesize}
\begin{sffamily}
\texttt{Johannes on 2022-08-03 at 23:10 said:}

Thank you, Cologero, for all the work you've put into this for all these years.

\hfill

\texttt{Rui Artur on 2022-08-04 at 08:16 said:}

This is incredibly useful! Thank you for this work.

\hfill

\texttt{Michael DeJoy on 2022-08-09 at 21:04 said:}

Cologero, the references to Matgioi's I Ching fill a void which I have been aware of for many years. Many interpretations of the hexagrams often are mechanical or try to appeal to Western analysis, especially Wilhelm with Jung's introduction. Some offer literal translations which try to lead the reader in a direction the author thinks they should go, so the I Ching will ``work'' for them. The translation you provided of Matgioi's work is the kind of explanation at the heart of the matter, it provides insight for the reader at the same time calms him down during a time the reader is looking for needed direction. Thank you, Michael

\hfill

\texttt{Rui Artur on 2022-08-12 at 05:26 said:}

It doesn't map exactly but there's something about this progression that seems to retrace Christ's steps from obscurity to glory. It also seems that Christ in a way breaks the `rules' of each, he is able to escape the `negatives' of each stage and hence progress (or perhaps progress despite the negatives).

Hidden Dragon -- the time before His ministry starts.

Dragon in the field -- `mine hour is not yet come' -- yet he changes the water into wine. these are the lesser mysteries, following the example of previous Mages.

Visible Dragon -- cleansing of the temple, healing, etc -- attracting the righteous, arousing the hatred of the Pharisees.

Jumping Dragon -- miracles of circumstance, Christ doesn't seek to do wonders for wonders sake. He helps those who appear in need before him (the paralytic, the 5000 who need food, the blind man).

Flying Dragon -- walking on water, the building of the Church, the `adoption' of St. John.

Hovering Dragon -- the Crucifixion, Resurrection and Ascension. Once the work is accomplished, He retreats.

Only a sketch, but I think there is something to it.

\hfill

\texttt{Luis on 2022-08-19 at 06:01 said:}

Thanks. This is very helpfull


\hfill

\end{sffamily}
\end{footnotesize}

